% Propietario conservado: /Users/ruben/Documents/ChatGPT/jueces y controles/REVISION_ARGUMENTAL_APERTURAS_CIERRES_20260915/delivery/ARTICLE_V_ES_ARGUMENTAL_20260915/manuscrito/sections/16_dinamica_temporal_finita.tex
% Extracto íntegro por residencia del cierre físico-matemático sucesor.
\section{Cristal temporal finito: excitación periódica, respuesta subarmónica y estabilidad espectral}
\label{v:sec:hmtf2-cristal-temporal}

\subsection{Excitación periódica inducida por el reloj de 108 estados}

Sea \(N=108\),
\[
 \mathcal H_{\rm clk}=\mathbb C^N,\qquad
 S|j\rangle=|j+1\rangle,\qquad
 Z|j\rangle=\zeta^j|j\rangle,\qquad
 \zeta=\exp(2\pi{\rm i}/N).
\]
En la base de Fourier \(\{|\widetilde k\rangle\}_{k=0}^{N-1}\), fijamos
\[
 H_{\rm clk}
 =\hbar_{\rm int}\frac{2\pi}{Nt_0}
 \sum_{k=0}^{N-1}k|\widetilde k\rangle\langle\widetilde k|,
 \qquad
 \exp\!\left(-\frac{{\rm i}}{\hbar_{\rm int}}H_{\rm clk}t_0\right)=S.
\]
Esta igualdad fija el reloj de partida. Para exhibir un impulso periódico no
constante, sea \(0<\eta<1\), \(T_{\rm d}=t_0\), y
definamos, en cada periodo,
\[
 H_{\rm d}(t)=
 \begin{cases}
  0,&0\leq t<\eta T_{\rm d},\\[2mm]
  (1-\eta)^{-1}H_{\rm clk},
  &\eta T_{\rm d}\leq t<T_{\rm d}.
 \end{cases}
 \qquad H_{\rm d}(t+T_{\rm d})=H_{\rm d}(t).
\]
Como los dos tramos conmutan, su propagador de Floquet es
\[
 U_F
 =\mathcal T\exp\!\left(
 -\frac{{\rm i}}{\hbar_{\rm int}}\int_0^{T_{\rm d}}H_{\rm d}(t)\,\mathrm dt
 \right)
 =\exp\!\left(-\frac{{\rm i}}{\hbar_{\rm int}}H_{\rm clk}T_{\rm d}\right)
 =S.
\]

Definimos los observables hermíticos
\[
 Q=\frac{Z+Z^\dagger}{2},
 \qquad
 P=\frac{Z-Z^\dagger}{2{\rm i}}.
\]

\begin{theorem}[Subarmonía primitiva del cristal temporal finito HMT]
\label{v:thm:hmtf2-subarmonia-108}
Para \(\rho_{j_0}=|j_0\rangle\langle j_0|\), el vector de orden
\[
 \mathbf q_n
 =\left(
 \operatorname{Tr}\rho_{j_0}U_F^{-n}QU_F^n,\,
 \operatorname{Tr}\rho_{j_0}U_F^{-n}PU_F^n
 \right)
\]
satisface
\[
 \mathbf q_n
 =\left(
 \cos\frac{2\pi(j_0+n)}N,\,
 \sin\frac{2\pi(j_0+n)}N
 \right).
\]
Su periodo primitivo es \(N T_{\rm d}=108T_{\rm d}\), mientras el drive
posee periodo \(T_{\rm d}\). Por tanto, el sistema realiza una respuesta
subarmónica primitiva de orden \(108\) en el dominio finito declarado.
\end{theorem}

\begin{proof}
La relación de Weyl \(S^{-1} Z S=\zeta Z\) da
\(S^{-n} Z S^n=\zeta^n Z\). Su parte real y su parte imaginaria producen la
fórmula. Si un entero \(d>0\) fuera un periodo, se tendría
\(\zeta^d=1\); la primitividad de \(\zeta\) obliga a \(N\mid d\).
\end{proof}

\subsection{Estabilidad espectral finita}

La separación mínima entre raíces distintas de la unidad es
\[
 g_N=\min_{k\ne\ell}|\zeta^k-\zeta^\ell|
 =2\sin\frac{\pi}{N}>0.
\]

\begin{theorem}[Rigidez covariante y estabilidad espectral finita]
\label{v:thm:hmtf2-estabilidad-cristal}
Se cumplen las dos afirmaciones siguientes.

\begin{enumerate}
 \item Para toda unitaria \(V\), la terna
 \[
  U_F^{(V)}=V U_F V^\dagger,\qquad
  Z^{(V)}=V Z V^\dagger,\qquad
  \rho^{(V)}=V\rho V^\dagger
 \]
 conserva exactamente la relación de Weyl, la subarmonía de orden \(N\) y
 la amplitud del vector de orden.
 \item Si \(U'\) es unitario y
 \(\|U'-U_F\|=\varepsilon<g_N/2\), los \(N\) discos de radio
 \(g_N/2\) centrados en \(\{\zeta^k\}\) separan de modo único los sectores
 espectrales de \(U'\). Sus proyectores de Riesz tienen el mismo rango que
 los de \(U_F\) a lo largo de toda deformación que permanezca en esa bola.
 Además, para todo observable \(A\), todo estado \(\rho\) y todo \(n\geq0\),
 \[
 \left|
 \operatorname{Tr}\rho\left(U'^{-n}AU'^n-U_F^{-n}AU_F^n\right)
 \right|
 \leq 2n\varepsilon\|A\|.
 \]
\end{enumerate}
\end{theorem}

\begin{proof}
La primera parte se obtiene conjugando cada identidad. Para la segunda, la
separación \(g_N\) impide que dos sectores entren en el mismo contorno; la
continuidad de los proyectores de Riesz mantiene sus rangos. La identidad
telescópica da
\(\|U'^n-U_F^n\|\leq n\varepsilon\), y la misma cota vale para las
potencias inversas. La desigualdad de expectativas resulta de insertar y
sustraer \(U'^{-n}AU_F^n\).
\end{proof}

El objeto obtenido es un \emph{cristal temporal finito HMT} con drive,
observable subarmónico y estabilidad espectral cuantificada. Que un material
macroscópico realice esta clase, conserve coherencia a gran volumen o
presente protección pretermal o localizada pertenece a
\textsc{reconocimiento externo}; no es un residual matemático de la
construcción finita.


